% C034 数学摘要；完整教学正文见正式六段文件。
直线 $Ax+By+C=0$，$A^2+B^2>0$；$a,b>0$。
消去 $y$ 必须 $B\ne0$；消去 $x$ 必须 $A\ne0$。
横开口 $x^2/a^2-y^2/b^2=1$，令 $M=B^2b^2-A^2a^2$、$H=C^2+M$。
\[
 Mx^2-2ACa^2x-a^2(C^2+B^2b^2)=0,\qquad
 My^2+2BCb^2y+b^2(C^2-A^2a^2)=0.
\]
$M\ne0$ 时，在合法消元方向用韦达：
\[
 x_1+x_2=\frac{2ACa^2}{M},\quad x_1x_2=-\frac{a^2(C^2+B^2b^2)}M;
\]
\[
 y_1+y_2=-\frac{2BCb^2}{M},\quad y_1y_2=\frac{b^2(C^2-A^2a^2)}M.
\]
纵开口 $y^2/a^2-x^2/b^2=1$，令 $N=A^2b^2-B^2a^2$、$K=C^2+N$。
\[
 Nx^2+2ACb^2x+b^2(C^2-B^2a^2)=0,\qquad
 Ny^2-2BCa^2y-a^2(C^2+A^2b^2)=0.
\]
$N\ne0$ 时对应韦达：
\[
 x_1+x_2=-\frac{2ACb^2}N,\quad x_1x_2=\frac{b^2(C^2-B^2a^2)}N;
\]
\[
 y_1+y_2=\frac{2BCa^2}N,\quad y_1y_2=-\frac{a^2(C^2+A^2b^2)}N.
\]
横开口 $M\ne0$ 时按 $H$ 正、零、负判断两交点、相切、无实交点；纵开口用 $(N,K)$。
$M=0$（或 $N=0$）时：$C\ne0$ 恰有一个有限交点且不是切线，$C=0$ 为渐近线且无有限交点。
横开口实弦须 $M\ne0,H>0$，有
\[
 S=\left(\frac{ACa^2}M,-\frac{BCb^2}M\right),\qquad PQ=\frac{2ab\sqrt{(A^2+B^2)H}}{|M|}.
\]
纵开口实弦须 $N\ne0,K>0$，有
\[
 S=\left(-\frac{ACb^2}N,\frac{BCa^2}N\right),\qquad PQ=\frac{2ab\sqrt{(A^2+B^2)K}}{|N|}.
\]
相切时按重根计韦达，长度为零；无实交点时不解释实弦。
